\documentclass[11pt,a4paper]{article}
\usepackage[utf8]{inputenc}
\usepackage[T1]{fontenc}
\usepackage[spanish]{babel}
\usepackage{amsmath,amssymb}
\usepackage{geometry}
\usepackage{tikz}
\usepackage{booktabs}
\usepackage{listings}
\geometry{margin=2.5cm}

\newcommand{\deter}[4]{\begin{vmatrix} #1 & #2 \\ #3 & #4 \end{vmatrix}}

\title{\'Area de pol\'igonos irregulares\\ F\'ormula de Gauss (``cord\'on de zapato'') con determinantes $2\times2$}
\author{Computaci\'on Gr\'afica}
\date{}

\begin{document}
\maketitle

\section{El problema}
Una parcela tiene un contorno irregular dado por sus v\'ertices ordenados
$P_1(x_1,y_1),\,P_2(x_2,y_2),\dots,P_n(x_n,y_n)$. Queremos su superficie
\textbf{sin} dividirla a mano en figuras simples, y que funcione igual con
$5$, $100$ o $1\,000$ v\'ertices (pol\'igonos simples, convexos o c\'oncavos).

\section{El determinante $2\times 2$ (Clase 04)}
Para dos v\'ertices consecutivos $P_i$ y $P_{i+1}$:
\[
D_i=\begin{vmatrix} x_i & x_{i+1} \\ y_i & y_{i+1} \end{vmatrix}
= x_i\,y_{i+1}-x_{i+1}\,y_i .
\]
(``Se multiplica en cruz y se resta''.)

\section{La f\'ormula de Gauss}
\begin{enumerate}
  \item Calcular $D_i$ para cada par consecutivo.
  \item Sumar todos, uniendo el \'ultimo v\'ertice con el primero ($P_{n+1}=P_1$):
  \[
  S=\sum_{i=1}^{n} D_i=\sum_{i=1}^{n}\left(x_i\,y_{i+1}-x_{i+1}\,y_i\right).
  \]
  \item El \'area es
  \[
  \boxed{A=\frac{|S|}{2}}
  \]
\end{enumerate}

\section{¿Por qu\'e funciona? (explicaci\'on sencilla)}
Tomemos el origen $O=(0,0)$ y el tri\'angulo $O\,P_i\,P_{i+1}$. Sus lados desde $O$
son los vectores $\vec{u}=(x_i,y_i)$ y $\vec{v}=(x_{i+1},y_{i+1})$. El
paralelogramo que forman tiene \'area $|D_i|$, as\'i que el tri\'angulo mide
$|D_i|/2$. Adem\'as, $D_i$ tiene \textbf{signo}:
\[
\frac{D_i}{2}=
\begin{cases}
+\text{\'area del tri\'angulo} & \text{si } P_i\to P_{i+1} \text{ gira en sentido antihorario respecto a } O,\\[2pt]
-\text{\'area del tri\'angulo} & \text{si gira en sentido horario}.
\end{cases}
\]
Al recorrer todo el contorno sumamos tri\'angulos con signo:
\begin{itemize}
  \item Las zonas que quedan \textbf{dentro} del pol\'igono se cuentan $+1$ vez neta.
  \item Las zonas que quedan \textbf{fuera} (entre el origen y el contorno) se suman
        con $+$ en un lado y con $-$ al regresar, as\'i que \textbf{se cancelan}.
\end{itemize}
Queda \'unicamente el \'area del pol\'igono, sin importar d\'onde est\'e el origen:
\[
\frac{S}{2}=\sum_{i=1}^{n}\frac{D_i}{2}=\pm A .
\]
El valor absoluto elimina el signo que depende del sentido de recorrido.

\section{Sentido del recorrido}
\[
S>0\Rightarrow \text{antihorario},\qquad
S<0\Rightarrow \text{horario},\qquad
S=0\Rightarrow \text{\'area nula (o contorno degenerado/autointersecante)}.
\]

\section{Ejemplos num\'ericos}

\subsection{Cuadrado de lado 4}
V\'ertices (antihorario): $P_1(0,0),\;P_2(4,0),\;P_3(4,4),\;P_4(0,4)$.
\begin{align*}
D_1&=\deter{0}{4}{0}{0}=0\cdot0-4\cdot0=0\\
D_2&=\deter{4}{4}{0}{4}=4\cdot4-4\cdot0=16\\
D_3&=\deter{4}{0}{4}{4}=4\cdot4-0\cdot4=16\\
D_4&=\deter{0}{0}{4}{0}=0\cdot0-0\cdot4=0\quad(P_4\to P_1)
\end{align*}
\[
S=0+16+16+0=32,\qquad A=\frac{|32|}{2}=\boxed{16}\ \text{u}^2,\qquad S>0\ \Rightarrow\ \text{antihorario}.
\]
Si se recorre al rev\'es ($P_1,P_4,P_3,P_2$) se obtiene $S=-32$ (horario) y el mismo
$A=16$.

\subsection{Tri\'angulo}
$P_1(0,0),\;P_2(4,0),\;P_3(0,3)$.
\begin{align*}
D_1&=0\cdot0-4\cdot0=0\\
D_2&=4\cdot3-0\cdot0=12\\
D_3&=0\cdot0-0\cdot3=0\quad(P_3\to P_1)
\end{align*}
\[
S=12,\qquad A=\frac{12}{2}=\boxed{6}\ \text{u}^2\ \text{(coincide con } \tfrac{b\,h}{2}=\tfrac{4\cdot3}{2}\text{)}.
\]

\subsection{Pol\'igono c\'oncavo en ``L''}
$P_1(0,0),\,P_2(4,0),\,P_3(4,2),\,P_4(2,2),\,P_5(2,4),\,P_6(0,4)$.
\begin{center}
\begin{tabular}{@{}cccc@{}}
\toprule
$i$ & Par & $x_iy_{i+1}-x_{i+1}y_i$ & $D_i$\\
\midrule
1 & $P_1\to P_2$ & $0\cdot0-4\cdot0$ & $0$\\
2 & $P_2\to P_3$ & $4\cdot2-4\cdot0$ & $8$\\
3 & $P_3\to P_4$ & $4\cdot2-2\cdot2$ & $4$\\
4 & $P_4\to P_5$ & $2\cdot4-2\cdot2$ & $4$\\
5 & $P_5\to P_6$ & $2\cdot4-0\cdot4$ & $8$\\
6 & $P_6\to P_1$ & $0\cdot0-0\cdot4$ & $0$\\
\bottomrule
\end{tabular}
\end{center}
\[
S=24,\qquad A=\frac{24}{2}=\boxed{12}\ \text{u}^2
\quad(\text{comprobaci\'on: } 4\cdot4-2\cdot2=12).
\]

\subsection{Caso rechazado: pol\'igono autointersecante (``corbata'')}
$P_1(0,0),\,P_2(4,4),\,P_3(4,0),\,P_4(0,4)$:
\begin{align*}
D_1&=0\cdot4-4\cdot0=0, & D_2&=4\cdot0-4\cdot4=-16,\\
D_3&=4\cdot4-0\cdot0=16, & D_4&=0\cdot0-0\cdot4=0.
\end{align*}
$S=0\Rightarrow A=0$, \textbf{¡un resultado falso!} Los dos tri\'angulos de la
corbata tienen orientaciones opuestas y se cancelan. Por eso la f\'ormula
\textbf{solo vale para pol\'igonos simples} y hay que validar antes.

\section{Pseudoc\'odigo}
\begin{verbatim}
ENTRADA: lista de n vertices P[1..n]
1. SI n < 3            -> ERROR "minimo 3 vertices"
2. SI no esta cerrado  -> ERROR "poligono no cerrado"
3. SI hay autointerseccion -> ERROR "poligono invalido"
4. suma <- 0
5. PARA i DESDE 1 HASTA n:
       j <- (i mod n) + 1              // el ultimo se une con el primero
       det <- P[i].x * P[j].y - P[j].x * P[i].y
       suma <- suma + det
6. area <- |suma| / 2
7. SI suma > 0 -> sentido <- antihorario
   SI suma < 0 -> sentido <- horario
SALIDA: area, sentido
\end{verbatim}

\section{Detecci\'on de autointersecci\'on (orientaci\'on de tres puntos)}
Para tres puntos $A,B,C$ se calcula el producto cruz
\[
o(A,B,C)=(B_x-A_x)(C_y-A_y)-(B_y-A_y)(C_x-A_x).
\]
\[
o>0:\ \text{$C$ est\'a a la izquierda de }\overrightarrow{AB},\quad
o<0:\ \text{a la derecha},\quad
o=0:\ \text{alineados}.
\]
Dos segmentos $AB$ y $CD$ \textbf{se cruzan} si
\[
\operatorname{sgn} o(A,B,C)\neq\operatorname{sgn} o(A,B,D)
\quad\text{y}\quad
\operatorname{sgn} o(C,D,A)\neq\operatorname{sgn} o(C,D,B),
\]
es decir, $C$ y $D$ quedan en lados opuestos de $AB$ \emph{y} $A$ y $B$ en lados
opuestos de $CD$. Si alg\'un $o=0$, se verifica adem\'as que el punto alineado caiga
dentro del segmento (caso de tocarse).

Se comparan todos los pares de aristas \textbf{no consecutivas} (las vecinas comparten
un v\'ertice, y eso es normal); en un pol\'igono cerrado la \'ultima arista tambi\'en
es vecina de la primera. Costo: $O(n^2)$ comparaciones.

\section{Escalabilidad y prueba con 1\,000 v\'ertices}
El c\'alculo del \'area hace $n$ determinantes: complejidad $O(n)$.

Para un pol\'igono regular de $n$ lados inscrito en un c\'irculo de radio $r$,
con $P_k=(r\cos\theta_k,\,r\sin\theta_k)$, $\theta_k=\tfrac{2\pi k}{n}$:
\[
D_k=r^2\left(\cos\theta_k\sin\theta_{k+1}-\cos\theta_{k+1}\sin\theta_k\right)
=r^2\sin\!\left(\tfrac{2\pi}{n}\right)
\]
\[
A_n=\frac{n}{2}\,r^2\sin\!\left(\frac{2\pi}{n}\right)\xrightarrow[n\to\infty]{}\pi r^2 .
\]
Con $n=1000$ y $r=10$:
\[
A_{1000}=500\cdot100\cdot\sin(0.00628319)\approx 314.1572,\qquad
\pi r^2\approx 314.1593,
\]
\[
\text{error}\approx 0.0021 \;(\approx 0.0007\,\%).
\]

\section{Resumen de casos de prueba}
\begin{center}
\begin{tabular}{@{}lcc@{}}
\toprule
Caso & $S$ & Resultado\\
\midrule
Cuadrado lado 4 & $32$ & $A=16$, antihorario\\
Tri\'angulo $(0,0),(4,0),(0,3)$ & $12$ & $A=6$, antihorario\\
``L'' c\'oncava & $24$ & $A=12$, antihorario\\
Corbata (autointersecante) & $0$ & \textbf{Rechazado}\\
C\'irculo, $n=1000$, $r=10$ & --- & $A\approx314.157$\\
\bottomrule
\end{tabular}
\end{center}

\end{document}

